% Compatibility shims for a newer local hyperref with the installed LaTeX kernel.
\providecommand\IfDocumentMetadataT[1]{}
\providecommand\IfDocumentMetadataTF[2]{#2}
\documentclass[aspectratio=169,9pt]{beamer}
\NeedsTeXFormat{LaTeX2e}
% Shared Agrawal homework style (loaded with \input from either deck).

\RequirePackage{fontspec}
\RequirePackage{amsmath,amssymb,mathtools}
\RequirePackage{booktabs}
\RequirePackage{graphicx}
\RequirePackage{microtype}
\RequirePackage{ragged2e}
\RequirePackage{hyperref}
\RequirePackage{tikz}
\usetikzlibrary{positioning}

\IfFontExistsTF{Avenir Next}{\setsansfont{Avenir Next}}{\setsansfont{TeX Gyre Heros}}

% Palette sampled from the supplied reference screenshot.
\definecolor{AgNavy}{HTML}{0C2852}
\definecolor{AgRed}{HTML}{982A34}
\definecolor{AgGraphite}{HTML}{000000}
\definecolor{AgMuted}{HTML}{000000}
\definecolor{AgRedTitle}{HTML}{F4E8EA}
\definecolor{AgRedBody}{HTML}{F9F5F5}
\definecolor{AgBlueTitle}{HTML}{EBEEF1}
\definecolor{AgBlueBody}{HTML}{F5F6F8}
\definecolor{AgDivider}{HTML}{D9DEE4}

\usetheme{default}
\usefonttheme{professionalfonts}
\setbeamertemplate{navigation symbols}{}
\setbeamersize{text margin left=0.85cm,text margin right=0.85cm}

\setbeamercolor{background canvas}{bg=white}
\setbeamercolor{normal text}{fg=AgGraphite,bg=white}
\setbeamercolor{frametitle}{fg=white,bg=AgNavy}
\setbeamercolor{structure}{fg=AgNavy}
\setbeamercolor{alerted text}{fg=AgRed}
\setbeamercolor{ag red title}{fg=AgRed,bg=AgRedTitle}
\setbeamercolor{ag red body}{fg=AgGraphite,bg=AgRedBody}
\setbeamercolor{ag blue title}{fg=AgNavy,bg=AgBlueTitle}
\setbeamercolor{ag blue body}{fg=AgGraphite,bg=AgBlueBody}
\setbeamercolor{ag accent line}{bg=AgRed}
\setbeamercolor{ag divider line}{bg=AgDivider}

\setbeamerfont{frametitle}{size=\large,series=\bfseries}
\setbeamerfont{normal text}{size=\normalsize}
\setbeamerfont{block title}{size=\normalsize,series=\bfseries}
\setbeamerfont{footline}{size=\scriptsize}

\setbeamertemplate{frametitle}{%
  \nointerlineskip
  \begin{beamercolorbox}[wd=\paperwidth,ht=0.62cm,dp=0.23cm,leftskip=0.40cm]{frametitle}%
    \usebeamerfont{frametitle}\insertframetitle
  \end{beamercolorbox}%
  \nointerlineskip
  \hbox{%
    \begin{beamercolorbox}[wd=1.05cm,ht=0.018cm,dp=0cm]{ag accent line}\end{beamercolorbox}%
    \begin{beamercolorbox}[wd=\dimexpr\paperwidth-1.05cm\relax,ht=0.018cm,dp=0cm]{ag divider line}\end{beamercolorbox}%
  }%
}

\setbeamertemplate{footline}{}

\setbeamertemplate{itemize item}{\color{AgNavy}\raise0.12ex\hbox{\scriptsize$\blacktriangleright$}}
\setbeamertemplate{itemize subitem}{\color{AgRed}\raise0.08ex\hbox{\tiny$\bullet$}}
\setlength{\leftmargini}{1.05em}
\setlength{\leftmarginii}{1.0em}

\newcommand{\E}{\mathbb{E}}
\newcommand{\Var}{\operatorname{Var}}
\newcommand{\Cov}{\operatorname{Cov}}

\newcommand{\agredblock}[2]{%
  \begin{beamercolorbox}[wd=\linewidth,sep=0.11cm]{ag red title}\usebeamerfont{block title}#1\end{beamercolorbox}%
  \nointerlineskip
  \begin{beamercolorbox}[wd=\linewidth,sep=0.14cm]{ag red body}#2\end{beamercolorbox}%
}

\newcommand{\agblueblock}[2]{%
  \begin{beamercolorbox}[wd=\linewidth,sep=0.11cm]{ag blue title}\usebeamerfont{block title}#1\end{beamercolorbox}%
  \nointerlineskip
  \begin{beamercolorbox}[wd=\linewidth,sep=0.14cm]{ag blue body}#2\end{beamercolorbox}%
}

\newcommand{\agcompactredblock}[2]{%
  \begin{beamercolorbox}[wd=\linewidth,sep=0.08cm]{ag red title}\usebeamerfont{block title}#1\end{beamercolorbox}%
  \nointerlineskip
  \begin{beamercolorbox}[wd=\linewidth,sep=0.09cm]{ag red body}#2\end{beamercolorbox}%
}

\newcommand{\agcompactblueblock}[2]{%
  \begin{beamercolorbox}[wd=\linewidth,sep=0.08cm]{ag blue title}\usebeamerfont{block title}#1\end{beamercolorbox}%
  \nointerlineskip
  \begin{beamercolorbox}[wd=\linewidth,sep=0.09cm]{ag blue body}#2\end{beamercolorbox}%
}

% Roomier variants used by the short deck: wider internal gutters without
% changing the more compact geometry of the 26-frame presentation.
\newcommand{\agshortredblock}[2]{%
  \begin{beamercolorbox}[wd=\linewidth,sep=0.12cm,leftskip=0.08cm,rightskip=0.08cm]{ag red title}\usebeamerfont{block title}#1\end{beamercolorbox}%
  \nointerlineskip
  \begin{beamercolorbox}[wd=\linewidth,sep=0.16cm,leftskip=0.10cm,rightskip=0.10cm]{ag red body}#2\end{beamercolorbox}%
}

\newcommand{\agshortblueblock}[2]{%
  \begin{beamercolorbox}[wd=\linewidth,sep=0.12cm,leftskip=0.08cm,rightskip=0.08cm]{ag blue title}\usebeamerfont{block title}#1\end{beamercolorbox}%
  \nointerlineskip
  \begin{beamercolorbox}[wd=\linewidth,sep=0.16cm,leftskip=0.10cm,rightskip=0.10cm]{ag blue body}#2\end{beamercolorbox}%
}

\newcommand{\agshortcompactredblock}[2]{%
  \begin{beamercolorbox}[wd=\linewidth,sep=0.10cm,leftskip=0.08cm,rightskip=0.08cm]{ag red title}\usebeamerfont{block title}#1\end{beamercolorbox}%
  \nointerlineskip
  \begin{beamercolorbox}[wd=\linewidth,sep=0.12cm,leftskip=0.12cm,rightskip=0.12cm]{ag red body}#2\end{beamercolorbox}%
}

\newcommand{\agshortcompactblueblock}[2]{%
  \begin{beamercolorbox}[wd=\linewidth,sep=0.10cm,leftskip=0.08cm,rightskip=0.08cm]{ag blue title}\usebeamerfont{block title}#1\end{beamercolorbox}%
  \nointerlineskip
  \begin{beamercolorbox}[wd=\linewidth,sep=0.12cm,leftskip=0.12cm,rightskip=0.12cm]{ag blue body}#2\end{beamercolorbox}%
}

\newcommand{\agshorttwocol}[4]{%
  \noindent\hspace*{0.18cm}%
  \begin{minipage}[t]{#1\linewidth}\vspace{0pt}#2\end{minipage}%
  \hfill
  \begin{minipage}[t]{#3\linewidth}\vspace{0pt}#4\end{minipage}%
  \hspace*{0.18cm}\par
}

\newcommand{\agtitlepage}[3]{%
  \setbeamertemplate{footline}{}%
  \vspace*{1.68cm}
  {\fontsize{17}{20}\selectfont\bfseries The Economics of Bicycles for the Mind\par}
  \vspace{0.18cm}
  {\fontsize{12.5}{15}\selectfont #1\par}
  \vspace{0.48cm}
  {\color{AgRed}\rule{\linewidth}{0.48pt}}
  \vspace{0.55cm}
  {\normalsize\bfseries Gabriel Saco\par}
  \vspace{0.06cm}
  {\normalsize August 24, 2026\par}
  \vspace{0.22cm}
  {\small Universidad del Pac\'ifico\par}
  \vspace{0.22cm}
  {\small\color{AgMuted}Agrawal, Gans \& Goldfarb (2025) \textperiodcentered\ NBER WP 34034\par}
  \vspace{0.12cm}
  {\small\href{https://github.com/gsaco/ai-02-agrawal}{\color{AgNavy}\ttfamily github.com/gsaco/ai-02-agrawal}\par}
  \vfill
  {\scriptsize\color{AgMuted}#2\hfill #3\par}
}

\newcommand{\agrestorefootline}{%
  \setbeamertemplate{footline}{}%
}

\setbeamersize{text margin left=0.55cm,text margin right=0.55cm}

\begin{document}

% 1/3 — title
\begin{frame}[plain]
  \agtitlepage
    {Cognitive tools, human judgment, and inequality}
    {}
    {AI Economic Modeling}
\end{frame}
\agrestorefootline

% 2/3 — paper and problem
\begin{frame}[t]{The paper and the agent's problem}
  \vspace{0.44cm}

  \agshortredblock{Question and single mechanism}{%
    \small
    \textbf{Question.} When a cognitive tool makes implementation easier, how do
    effort, productivity, returns to human abilities, and inequality change?

    \vspace{0.08cm}
    \textbf{Mechanism.} The tool is a \alert{substitute for implementation effort}
    at the margin; implementation skill $s$, payoff judgment $\alpha$, and
    opportunity judgment $\gamma(t)$ determine who can exploit it.
  }

  \vspace{0.40cm}

  \agshortblueblock{Agent's problem, conditional on finding an opportunity}{%
    \agshorttwocol{.535}{%
      \small Choose effort $e_t$ to maximize expected improvement value net of cost:
      \[
        \boxed{e_t^*(\theta)\in\arg\max_{e_t\ge0}
        \{p(se_t;\theta)\alpha\Delta-c(e_t;\theta)\}.}
      \]
      \vspace{-0.18cm}
      {\small Interior FOC:
      $p'(se_t^*;\theta)s\alpha\Delta=c'(e_t^*;\theta)$.}
    }{.385}{%
      \small
      \textbf{Choice and constraint:} $e_t\ge0$.\\[0.10cm]
      \textbf{Technology:} $p\in[0,1]$ is increasing and weakly concave.\\[0.10cm]
      \textbf{Cost:} $c$ is increasing and weakly convex.\\[0.10cm]
      \textbf{Tool:} raises $p$ and/or lowers $c$, while lowering $p'/c'$.
    }
  }
\end{frame}

% 3/3 — corrected exact result
\begin{frame}[t]{Proposition 3: the conditional variance result}
  \setlength{\abovedisplayskip}{0.32em}
  \setlength{\belowdisplayskip}{0.32em}
  \vspace{0.10cm}

  \agshortcompactblueblock{Specialization and conditions}{%
    \small
    \agshorttwocol{.445}{%
      $p(se;\theta)=\sqrt{se+\theta}$, $c(e)=e$;
      $\Gamma=\gamma_0/(1-\delta\gamma)$ and $\delta\gamma<1$.
    }{.445}{%
      $\alpha,\gamma_0,\gamma,s$ are independent, have positive support,
      and satisfy $\mu_i>3\sigma_i$.
    }
    \vspace{0.07cm}
    Common interior range:
    $\theta<\alpha^2\Delta^2s^2/4$ for every worker under comparison.
  }

  \vspace{0.18cm}

  \agshortcompactredblock{Corrected exact statement}{%
    \small
    \agshorttwocol{.445}{%
      \[
      V(\theta)=\Gamma\left(\frac{\alpha^2\Delta^2s}{4}+\frac{\theta}{s}\right).
      \]
      \[
      \frac{d\E[V]}{d\theta}=\E[\Gamma]\E[1/s]>0.
      \]
      The U-shape concerns cross-sectional $\Var[V(\theta)]$ as continuous
      tool quality $\theta$ changes.
    }{.445}{%
      If
      \[
      \frac{\E[\Gamma^2]}{\E[\Gamma]^2}<\mu_s\E[1/s],
      \qquad \Var(\Gamma/s)>0,
      \]
      then $\Var[V(\theta)]$ is strictly convex, falls for
      $\theta<\theta^*$, and rises for $\theta>\theta^*$, where
      $\theta^*=-a_0/[2\Var(\Gamma/s)]>0$.
    }
  }

  \vspace{0.18cm}

  \agshortcompactblueblock{Endpoint correction and intuition}{%
    \small
    Equation (32)'s positive slope at $\theta=1$ also requires
    \alert{$\theta^*<1$}. \textbf{Intuition:} inverse skill bias equalizes first;
    heterogeneous opportunity judgment can dominate later.
  }
\end{frame}

\end{document}
